\chapter{CONCLUSÃO} \label{chap_conclusao}

Neste trabalho, buscamos investigar o desempenho atual de um computador quântico na tarefa de classificação de dados da base de dados Iris, comparando-o com um classificador clássico. O objetivo principal foi avaliar a capacidade do computador quântico real em relação a um problema de aprendizado de máquina, especificamente na área de classificação de dados. A escolha da base de dados Iris como caso de estudo se deve à sua ampla utilização na comunidade científica como benchmark para algoritmos de classificação.

Para alcançar esse objetivo, utilizamos um método de aprendizado de máquina quântico supervisionado, o classificador quântico variacional, executado em um computador quântico. O classificador quântico foi construído a partir de circuitos quânticos para codificação de informações clássicas em estados quânticos e para aprendizagem. O treinamento consistiu em aprender os parâmetros das portas quânticas que otimizam a função de custo escolhida.

Ao realizar comparações detalhadas entre o classificador clássico e o computador quântico real, analisamos diversos aspectos, como métricas de desempenho e tamanhos de conjunto de treinamento. Observamos que, de maneira geral, o classificador clássico demonstrou um desempenho superior ao quântico para a maioria das métricas em todos os tamanhos de treinamento avaliados.

Esses resultados sugerem que, para o conjunto de dados Iris utilizado neste estudo, o modelo clássico pode ser mais adequado. No entanto, é importante ressaltar que o classificador quântico exibiu um desempenho consistente em todas as métricas e tamanhos de treinamento.

Apesar do classificador quântico ter apresentado resultados inferiores aos do classificador clássico, é fundamental considerar que a computação quântica ainda é um campo emergente. Levando em conta a consistência e a robustez demonstradas pelo classificador quântico neste estudo, espera-se que, com mais otimizações e desenvolvimentos nessa área, seja possível alcançar e até mesmo superar o desempenho do classificador clássico.

Este trabalho contribuiu significativamente para a compreensão do estado atual da computação quântica em relação ao aprendizado de máquina. Ao analisar a capacidade de um computador quântico real na tarefa de classificação de dados, foram identificadas suas limitações e obtidos \textit{insights} valiosos para futuras pesquisas e desenvolvimentos nessa área em constante evolução.

Além disso, este estudo abre caminho para futuras pesquisas, particularmente aquelas voltadas para a otimização dos classificadores quânticos e para a exploração de suas vantagens potenciais. A expansão dessa análise para outros conjuntos de dados e cenários pode fornecer insights adicionais sobre a aplicabilidade e eficácia dos classificadores quânticos.